\chapter{复制：细胞怎样抄写几十亿个字母}

\begin{kaichang}
切菜时不小心在手指上划了一道小口子。你冲了冲水，贴上创可贴，几天后揭下来一看：痂掉了，底下长出了新皮肤，连指纹的纹路都和原来几乎一样。

现在请你猜两件事。第一，为了补好这道小口子，你的身体新造了多少个细胞？第二，每一个新细胞里都装着一整套完整的遗传说明书——人类的这套说明书有三十亿个“字母”（碱基），这些字母是从哪儿来的？是新细胞自己现编的吗？还是从旧细胞那里抄来的？如果是抄的，抄写员是谁，怎么保证不抄错？

先写下你的猜测。这一章结束时，我们会回到这道小伤口，看看它底下刚刚发生过一场多么浩大的抄写工程。
\end{kaichang}

\section{先翻一倍，再分家}

伤口能长好，是因为伤口边缘的细胞在分裂：一个变两个，两个变四个，把缺口填满。细胞分裂这件事你在第 64 章已经见过——染色体先复制，再被均分到两个子细胞里。这里我们要把镜头推进到分子层，盯住其中一个环节：染色体里的 DNA，是怎样“先翻一倍”的？

先看几个能观察、能测量的事实。你的指甲每周长出一毫米左右，头发每个月长一厘米多，这些新长出来的部分全是新细胞；培养皿里一个大肠杆菌在 $37\,^{\circ}\mathrm{C}$ 下每二十来分钟就变成两个；一个受精卵经过几十轮分裂，变成你身体里几十万亿个细胞。所有这些分裂有一个共同的硬规定：\textbf{分裂之前，细胞里的 DNA 总量必须先恰好翻一倍}。这一点是可以直接测出来的——给 DNA 染上会发荧光的染料，在显微镜下量亮度，正准备分裂的细胞亮度正好是普通细胞的两倍，不多也不少。

为什么必须是“恰好一倍”？多了少了会怎样？想想就明白：每个子细胞都必须拿到一整套说明书，一页不能多，一页不能少。少抄了一段，子细胞就缺了某个零件的图纸；多抄了一段，遗传信息就乱了套。所以细胞把“抄书”和“分家”安排成两个先后分明的阶段：先把三十亿个字母完整抄一遍，再把自己一分为二。

问题是：三十亿个字母，怎么抄？靠谁来抄？抄错了怎么办？要回答这些，得回到第 66 章那张双螺旋的图纸。

\section{一条链，就是另一条的图纸}

回忆双螺旋的结构：两条长链互相缠绕，中间靠碱基配对咬合——腺嘌呤（A）永远对着胸腺嘧啶（T），鸟嘌呤（G）永远对着胞嘧啶（C）。第 66 章说过，这种配对有严格的化学原因：碱基的形状和氢键供体受体的位置决定了只有这两对能严丝合缝地卡进双螺旋。

现在请注意一个看似平常、实则惊人的推论：\textbf{只要知道一条链的序列，另一条链的序列就被完全决定了}。一条链上写 A，对面那条必然是 T；写 G，对面必然是 C。两条链互为底片和照片——拿着任何一条，都能把另一条重新洗出来。

1953 年沃森和克里克在论文结尾写了一句后来被反复引用的话：“我们注意到，我们假定的这种配对方式，直接暗示了遗传物质一种可能的复制机制。”他们指的就是这个：把两条链像拉链一样拉开，每条链各自当模板，按照配对规则在空位上逐个安上新的碱基——A 的位置放 T，G 的位置放 C。抄完之后，原来的一份双螺旋变成两份双螺旋，每一份都由一条旧链和一条新链组成。

等一下，“每一份都是一新一旧”——这是一定的吗？其实不一定。光看逻辑，至少有三种抄法都说得通：

\begin{itemize}[itemsep=2pt]
\item \textbf{半保留复制}：两条旧链拆开，各带一条新链。抄完的两份 DNA，每份都是“一半旧、一半新”。
\item \textbf{全保留复制}：旧双螺旋原封不动地待在一起，另外凭空抄出一份全新的双螺旋。抄完一份全新、一份全旧。
\item \textbf{分散复制}：旧链被打成碎片，和新合成的材料混着拼接。两份 DNA 的每条链都是新旧相间的“百衲衣”。
\end{itemize}

三种方案都能让 DNA 翻倍，也都符合“互补配对”的规则。1950 年代的生物学家为此争论了好几年，谁也说服不了谁——因为这恰恰是光靠想、想不出来的问题。你得让实验来投票。这场投票我们放到后面的证据层细讲，先把抄写过程本身看清楚。

\begin{shiyan}
\textbf{把 DNA 亲手捞出来。}DNA 不是课本上的抽象符号，它是真实的长分子，长到你能用肉眼看见——只要一次捞出来足够多条。

材料：两三颗草莓（或半根香蕉）、一小勺食盐、几滴洗洁精、水、纱布或咖啡滤纸、两个透明杯子、一小杯冰镇的高度酒精（如伏特加或白酒，请家长提供）。

步骤：草莓去蒂装进保鲜袋，加两勺水、一小撮盐和几滴洗洁精，隔着袋子轻轻捏碎成糊，静置五分钟；用纱布把果糊滤进杯子，留小半杯滤液；把冰酒精沿杯壁缓缓倒入，让它浮在滤液上，不要搅拌；静置观察两层交界处。

观察点：几分钟后，交界处会出现白色絮状物，可以用牙签卷起来。那就是成千上万条缠在一起的 DNA（混着些蛋白质）。盐和洗洁精负责拆散细胞膜和蛋白质，DNA 不溶于冰冷酒精，于是抱团析出。想想看：一颗草莓的细胞里，这些“絮”原本折叠在比针尖还小的空间里。

安全提示：酒精易燃，全程远离火源；实验材料处理后不可食用；结束后洗手。本实验可以在家做，但酒精取用请由家长操作。
\end{shiyan}

\section{一支分工明确的抄写队}

配对规则解决了“抄什么”，但真正的抄写是酶干的活。细胞内没有抄写员的桌子，有的是一台台蛋白质机器，沿着 DNA 跑动，像一支流水线上的施工队。认识一下这支队伍的主要成员。

\textbf{解旋酶}是开路先锋。它卡在双螺旋上，消耗 ATP（第 52 章的能量货币），硬生生把两条链之间的氢键扯开，在前方撕出一个 Y 形的开口——这个开口叫\textbf{复制叉}，就是抄写作业的工地。别忘了，两条链是互相缠绕的螺旋，前面不断拆开，更前方就会越拧越紧，像拧毛巾一样拧出疙瘩，所以还有一种酶负责在更前方把“拧劲”卸掉。

\textbf{引物酶}负责垫第一块砖。这里有个很关键的规矩：负责抄写的 DNA 聚合酶有个怪脾气——它\textbf{不能从头开始}，只能把新碱基接到一段已经存在的链的末端上。所以每段抄写开工前，引物酶先合成一小段 RNA 当“起跑踏板”，后面的碱基全接在这段踏板上。踏板用完即拆，缺口再被补上 DNA。

\textbf{DNA 聚合酶}是主力抄写员。它沿着模板链滑动，逐个读碱基，从细胞质里的原料库（四种脱氧核苷酸，各自拖着三个磷酸基团）中挑出能配对的那个，咔哒一声接到新链的末端。接上去所需的能量，就来自原料自身甩掉两个磷酸基团时释放的化学能——注意，是成键和水解反应整体放能，不是“某根高能键里藏着能量”（第 52 章澄清过这个误解）。

\textbf{连接酶}是收尾的泥瓦匠，负责把新链上的缝隙焊死。为什么需要它？这就牵出复制过程里最拧巴的一个设计。

聚合酶还有个改不掉的方向癖：它只朝一个方向干活。而 DNA 的两条链是反向平行的（第 66 章），复制叉又在不断向前撕开。结果是：以其中一条链为模板，新链可以跟着叉子一路连续抄下去，这条叫\textbf{前导链}；以另一条链为模板，聚合酶只能背对着叉子前进方向、一小段一小段地“倒车施工”，每抄完一小段，就得跳回叉子附近重新垫引物、再抄下一段。这条断断续续的链叫\textbf{后随链}，那些小片段叫冈崎片段——以发现它的日本科学家冈崎令治夫妇命名。每个冈崎片段的头上有引物要拆、片段之间有缝要焊，所以后随链的活儿格外繁琐，连接酶主要就是在这儿忙活。

\begin{qiaoliang}
这支队伍干活有多快？细菌里，一个复制叉每秒能接上约 1000 个碱基。大肠杆菌的染色体约 460 万个碱基对，从一个起点朝两个方向各派一个复制叉，两路夹攻，抄完整条染色体需要 $4.6\times10^{6}\div(2\times1000)\approx2300$ 秒，也就是四十分钟上下——跟它二十分钟一代的分裂节奏基本吻合（它的诀窍是：上一轮还没抄完，下一轮已经在半成的副本上提前开工）。

人的复制叉慢得多，每秒只有约 50 个碱基。要是只靠一个起点双向抄，三十亿个碱基对得抄将近两年。可你的细胞实际只用大约 8 小时就抄完了。怎么办到的？答案是\textbf{多点开工}：人类染色体上散布着几万个复制起点，细胞分批点燃它们，同一时刻有成千上万个复制叉同时施工。慢不要紧，工地多就行。
\end{qiaoliang}

\section{抄错了怎么办：自带橡皮的铅笔}

速度惊人，准确率更惊人。先说原料端：聚合酶挑碱基主要靠配对的几何形状，但偶尔也会看走眼——大约每抄一万到十万个字母，就会塞进一个错的。如果错误率到此为止，你的基因组每复制一次就要多出几十万个拼写错误，后果不堪设想。

所以聚合酶随身带着橡皮：它手里有个“校读”工位，刚接上去的碱基如果配对不齐，新链末端的形状就别扭，聚合酶会卡住、倒退、把刚接的碱基切掉重来。这道校对把错误率压低到大约千万分之一。抄完之后，还有第二道防线：一套\textbf{错配修复}系统沿着新抄的 DNA 巡查，专找漏网的错配，认出哪条是新链（细胞有办法给新旧链做记号），把错的那段挖掉重抄。两关过后，最终错误率降到每抄十亿到一百亿个字母才错一个。

这是个什么概念？假如让你手抄一本书，每抄十亿个字母才允许错一个——你得抄完两千多套《红楼梦》全本，才有资格错一个字。细胞做到了，靠的是分子层面的“写—查—修”三道工序，而不是运气。

但请注意：十亿分之一\textbf{不是零}。三十亿个碱基对抄一遍，平均下来每次细胞分裂总会留下零星的几个新拼写差异。你的每一个体细胞，基因组都和受精卵的原版有着细微出入。这个事实先记在心里，它是后面好几章的入口——本章末尾还会回来接它。

\section{把抄写作业写成符号}

讲了这么多机器，该把抄写作业本身落成符号了。规矩只有两条：配对（A 对 T，G 对 C）加方向（新链永远从 $5'$ 端往 $3'$ 端延长，模板则被反向读取）。比如一小段模板链：

\[3'-\mathrm{TACGGATC}-5'\]

抄写员从左往右读它，配出的新链就是：

\[5'-\mathrm{ATGCCTAG}-3'\]

一个字母对一个字母，两条链方向相反——这就是第 66 章“反向平行”在抄写时的实际后果，也是后随链不得不“倒车施工”的根本原因。你可以自己随便写一串字母当模板，按这两条规矩配出互补链：恭喜，你刚刚用纸和笔完成了一次 DNA 聚合酶的工作，只是慢了一亿倍。

别小看这张字母表的分量。整部遗传说明书只有四个字母，比英文的二十六个还寒酸，却靠“排列顺序”一项本事装下了建造一个完整的人所需的全部信息——第 66 章说过，信息不在字母本身，而在顺序里。复制要守护的正是顺序：机器可以慢、可以绕路，唯独字母的先后一个都不能乱。

\section{三种抄法，实验投票}

\begin{zhengju}
回到前面那个悬案：半保留、全保留、分散复制，到底哪种是真的？1958 年，美国加州理工学院的梅塞尔森和斯塔尔设计了一个被后世称为“生物学中最漂亮的实验”的方案，一锤定音。

他们的绝招是给 DNA“称体重”。氮是 DNA 碱基的重要成分，普通氮原子是 $\chem{^{14}N}$，还有一种稳定的重同位素 $\chem{^{15}N}$（第 7 章讲过，同位素质子数相同、中子数不同）。研究者先把大肠杆菌在只含 $\chem{^{15}N}$ 的培养液里养上十几代，让每条 DNA 链都变成“重链”；然后把细菌突然移到普通 $\chem{^{14}N}$ 培养液里，从此新合成的链都是“轻链”。每过一代取一次样，把 DNA 提取出来放进氯化铯溶液中超速离心十几个钟头——溶液在离心力下形成上轻下重的密度梯度，DNA 分子会漂在和自己密度恰好相等的那一层，重的沉得深，轻的浮得高，用紫外光一照就能看到一条条带的位置。

三种抄法给出的预言完全不同。全保留说：第一代就该分成两条带——一条全重（旧双螺旋），一条全轻（新双螺旋）。分散说：每一代都只有一条带，只是位置一代比一代往上浮。半保留说：第一代只有一条带，停在正中间（每条双螺旋一重一轻）；到第二代，这条中带一分为二，一半还在中间，另一半浮到轻的位置。

实验结果：第一代，只有中间一条带——全保留当场出局。第二代，中带和轻带各占一半——分散复制也出局。梅塞尔森和斯塔尔后来还加热把第二代的杂合 DNA 拆成单链再离心，果然拆出一条重链一条轻链，把“每条双螺旋一新一旧”钉到了最后的棺材板上。这个实验没有靠任何天才的灵光，靠的是把三种假设的\textbf{可观测后果}提前算清楚，再让离心管投票。
\end{zhengju}

\begin{tikzpicture}
\draw[->] (-0.6,0) -- (-0.6,3.2) node[left]{密度大};
\foreach \x/\lab in {0/第零代, 2.4/第一代, 4.8/第二代}{
  \draw (\x,0.3) -- (\x,2.7);
  \draw (\x+0.8,0.3) -- (\x+0.8,2.7);
  \draw (\x,0.3) -- (\x+0.8,0.3);
  \draw (\x,2.7) -- (\x+0.8,2.7);
  \node at (\x+0.4,-0.3){\lab};
}
\draw[line width=2pt] (0,0.7) -- (0.8,0.7);
\draw[line width=2pt] (2.4,1.6) -- (3.2,1.6);
\draw[line width=2pt] (4.8,1.6) -- (5.6,1.6);
\draw[line width=2pt] (4.8,2.5) -- (5.6,2.5);
\node at (6.6,0.7){全重};
\node at (6.6,1.6){一重一轻};
\node at (6.6,2.5){全轻};
\end{tikzpicture}

这是离心结果的示意图：横线是 DNA 带停留的位置，不代表真实照片。注意第二代同时出现中带和轻带——这一条就足以排除分散复制。

\section{队尾的麻烦：端粒}

这套抄写机器近乎完美，但它有一个天生的死角，藏在染色体的两头。

回忆后随链的施工方式：每抄一小段都要先垫 RNA 引物，事后拆掉引物再补 DNA。问题在于染色体最末端：最后一小段引物拆掉之后，它的\textbf{前面再没有空位}可以让聚合酶落脚补抄——聚合酶只能往后接，不能往前补。于是每复制一轮，后随链这一头就要短上一小截。这叫“末端复制难题”，是抄写机制的方向癖直接推导出来的必然结果，不是机器故障。

细胞的应对是给染色体两端各扣一顶“消耗型安全帽”：\textbf{端粒}，几千个无编码意义的重复序列（人的端粒是 TTAGGG 重复几千遍）。每次分裂磨掉的是帽子，不是正文的图纸。但帽子总有磨薄的一天。培养皿里的人类细胞分裂四五十次后，端粒短到警戒线，细胞就永久停止分裂——这个上限叫海弗利克极限。第 60 章讲稳态时说过身体在不停纠偏，端粒其实是纠偏系统的另一面：它给每个细胞的“分裂次数”设了额度，让无限增殖变得不可能。

那生殖细胞、干细胞凭什么代代相传？它们有一种专门的酶——\textbf{端粒酶}，自带一小段 RNA 模板，能反过来给端粒“续帽”。而绝大多数普通体细胞把它关了。令人警惕的是，约九成的癌细胞重新打开了端粒酶，把消耗型安全帽变成了永久牌，这是它们能无限分裂的门票之一。所以端粒酶既是生命延续的工具，也是癌症研究的重要靶点——同一把钥匙，开的是哪扇门，取决于握在谁手里。

顺带交代一下本章模型的边界。我们讲的是“细胞核里的染色体怎样复制”这条主线，细节是按主流教材简化的：真实细胞里干活的 DNA 聚合酶不止一种，有的专管开工，有的专管前导链，有的专管后随链，分工比“一个抄写员”精细得多；错误率的具体数字也因物种、组织、年龄而异，十亿到一百亿分之一只是一个数量级。细菌是环状染色体，没有末端，天然绕开了端粒的麻烦；许多病毒干脆用 RNA 当遗传物质，抄写规则和易错程度又是另一本账（第 65 章提过 RNA 病毒，第 79 章还会再遇到它们）。主线图景在这些场合依然好用，但别把它当成放之四海皆准的施工图。

\begin{wuqu}
“DNA 复制出错，都是因为辐射、毒物这些外界祸害；只要环境干净，复制就是绝对准确的。”——错。即使没有半点辐射和毒物，抄写机器本身就有万分之一的原料级失误，校对加修复之后仍有十亿分之一量级的漏网之鱼。三十亿个碱基对抄一遍，平均注定留下零星几个新差异——你身体里今天分裂的每一个细胞，几乎都和受精卵的原版不完全相同。况且错误的后果大多没你想的严重：绝大多数差异落在不编码的序列里，或者根本不影响蛋白质的工作。真正需要警惕的，是恰好落在关键基因上、又躲过修复的那极少数——突变究竟是错误、材料还是创新，第 72 章会专门掰扯。
\end{wuqu}

\section{你身体里的抄写总量}

来算一笔大账，体会一下这场抄写工程的规模。你的身体里每天大约有三千亿个细胞走到生命尽头并被替换（主要是血细胞和肠道内壁细胞），每个新细胞上岗前都要把整套基因组抄一遍。人类二倍体基因组约 $6\times10^{9}$ 个碱基对，于是你每天的抄写总量大约是：

\[3\times10^{11}\ \text{个细胞}\times 6\times10^{9}\ \text{碱基对}\approx 2\times10^{21}\ \text{个字母}\]

$2\times10^{21}$ 个字母是什么概念？假如把每个字母印成书上一个字符，这套书堆起来可以装满几百万座图书馆——而这一切在你体内悄无声息地发生，错误总数还被压在极低的水平。生命维持自身的日常开支，大得超出直觉。

医学还反过来借用了这套抄写机器。核酸检测和法医鉴定用的 PCR 技术（聚合酶链式反应），就是把解旋、配对、延长这三步搬进试管：加热到 $95\,^{\circ}\mathrm{C}$ 左右代替解旋酶把 DNA 撕开，降温让人工合成的引物贴上去，再用一种从温泉细菌里借来的耐高温聚合酶延长新链。每循环一次，目标片段翻一倍。

\begin{tiaozhan}
PCR 的翻倍是指数增长：从 1 份拷贝出发，第 $n$ 次循环后有 $2^{n}$ 份。10 次循环约一千份，20 次约一百万份，30 次约十亿份（$2^{30}\approx1.07\times10^{9}$）。这就是为什么咽拭子里哪怕只有几个病毒的一小段核酸，三十来个循环后也能被仪器检出——指数增长是微量信号最凶猛的放大器。反过来，它也解释了为什么 PCR 实验室最怕污染：一个不该出现的分子同样会被放大十亿倍。用对数算一算：要把 1 份放大到 $10^{12}$ 份，至少需要多少次循环？（提示：$\lg 2\approx0.30$。）跳过本框不影响主线。
\end{tiaozhan}

\section{回到那道小伤口}

现在可以揭开创可贴了。你那道小伤口愈合时，边缘的上皮细胞和成纤维细胞加班加点分裂，填平缺口。每个新细胞上岗前，都经历了一场完整的抄写：几万个起点同时开工，解旋酶撕开双螺旋，引物垫路，聚合酶以每秒约 50 个字母的速度推进，校对橡皮跟在后面擦，错配修复系统再巡一遍，最后连接酶焊好所有缝隙——几个小时，六十亿个字母，误差控制在个位数。指纹的纹路能和原来几乎一样，靠的不是运气，而是这套“写—查—修”流水线把信息几乎原封不动地传了下去。

但也别忘掉那“个位数”的误差，还有每次分裂都在磨损的端粒。抄写极其可靠，却不是永动机，更不是绝对的完美——细胞怎样处理抄错的内容，身体又怎样给分裂次数记账，都是生命故事里真正的悬念。

还有一个更根本的问题，我们从头到尾都绕开了：这套三十亿字母的说明书，每一页到底写了什么？细胞怎么把某一页的“图纸”变成真正能干活的蛋白质？抄写只是保存信息，\textbf{使用}信息是另一套流程——第一步是把某一页从锁在细胞核里的母本上转抄成一份可以带出门的工作单。下一章，我们就跟着这份工作单走。

\section*{练一练}

\begin{sikaoti}
\item 随意写一串 12 个字母的 DNA 序列（只用 A、T、G、C），标出 $5'$ 和 $3'$ 端，然后写出复制时以它为模板合成的新链，并标明新链的方向。同桌互换检查：配对和方向都符合规则吗？
\item 画一幅复制叉的粒子示意图：画出一个 Y 形开口，标出模板双链、解旋酶的位置、前导链、后随链和冈崎片段，用箭头标出每条新链的延长方向。图旁注明：这是帮助思考的表示法，不代表分子的真实外观。
\item 梅塞尔森—斯塔尔实验中，假如 DNA 采用的是分散复制，第一、二、三代离心管里各会看到怎样的条带？画三条“试管”表示，并和半保留复制的实际结果对照，指出哪一代的结果能一锤定音地区分两种假说。
\item 人的一个复制叉每秒抄约 50 个碱基，S 期约 8 小时。估算一个复制叉在整个 S 期大约能抄多少碱基对，再估算抄完三十亿碱基对的基因组，需要先后动用多少个复制起点（假设起点不重复使用）。列出你的计算过程。
\item 体细胞每分裂一次端粒就短一截，而癌细胞重新打开了端粒酶。有人据此宣称“补充端粒酶就能长生不老”。根据本章内容和剂量、证据原则，指出这个推理至少有两处跳跃，并说明可能的风险。
\item 为什么错配修复系统必须先分清“哪条是旧链、哪条是新链”才能纠错？如果它认错了链，把旧链当成错的改掉，会发生什么？用信息流图（模板→新链→修复）说明你的推理。
\end{sikaoti}
